\documentclass[letterpaper,11pt,twoside]{article}
\usepackage{graphicx} % Required for inserting images
\usepackage[table,xcdraw,dvipsnames]{xcolor}
\usepackage{amsmath,amsfonts,amssymb,amsthm}
\usepackage{listings}
\usepackage{lipsum}
\usepackage{hyperref}
\usepackage{enumitem}

\usepackage{tikz}
\usepackage[siunitx, RPvoltages]{circuitikz}
\usetikzlibrary{3d}
\usepackage{comment}
\usepackage{caption,subcaption}
\usepackage{pgfplots}
\usepackage{dsfont}
\pgfplotsset{compat=newest} % or a newer version if available
\usepgfplotslibrary{groupplots}
\usetikzlibrary{pgfplots.groupplots}
\usetikzlibrary{shapes.geometric, arrows}
\tikzstyle{arrow} = [->,>=stealth,shorten >=2pt]
\newcommand{\re}[1]{\text{Re}\left(#1\right)}
\newcommand{\im}[1]{\text{Im}\left(#1\right)}
\newcommand{\bA}{\bm{A}}
\newcommand{\bB}{\bm{B}}
\newcommand{\bC}{\bm{C}}
\newcommand{\bE}{\bm{E}}
\newcommand{\E}{\mathcal{E}}
\newcommand{\bH}{\bm{H}}
\newcommand{\bD}{\bm{D}}
\newcommand{\bO}{\bm{0}}
\newcommand{\br}{\bm{r}}
\newcommand{\bk}{\bm{k}}
\newcommand{\bR}{\bm{R}}
\newcommand{\hA}{\hat{A}}
\newcommand{\hB}{\hat{B}}
\newcommand{\hC}{\hat{C}}
\newcommand{\hU}{\hat{U}}
\newcommand{\ha}{\hat{a}}
\newcommand{\hb}{\hat{b}}
\newcommand{\hc}{\hat{c}}
\newcommand{\hn}{\hat{n}}
\newcommand{\hH}{\hat{H}}
\newcommand{\hN}{\hat{N}}
\newcommand{\hr}{\bm{\hat{r}}}
\newcommand{\hbE}{\hat{\bm{E}}}
\newcommand{\hx}{\hat{x}}
\newcommand{\hX}{\hat{X}}
\newcommand{\hI}{\hat{I}}
\newcommand{\hp}{\hat{p}}
\newcommand{\hP}{\hat{P}}
\newcommand{\hS}{\hat{S}}
\newcommand{\hD}{\hat{D}}
\newcommand{\hy}{\hat{y}}
\newcommand{\hrho}{\hat{\rho}}
\newcommand{\hz}{\hat{z}}
\newcommand{\laplace}[1]{\mathscr{L}\left[#1\right]}
\newcommand{\ilaplace}[1]{\mathscr{L}^{-1}\left[#1\right]}
\newcommand{\fourier}[1]{\mathscr{F}\left[#1\right]}
\newcommand{\Tr}[1]{\text{Tr}\left[#1\right]}
\newcommand{\ifourier}[1]{\mathscr{F}^{-1}\left[#1\right]}
\newcommand{\ket}[1]{\left|#1\right\rangle}
\newcommand{\bra}[1]{\left\langle#1\right|}
\newcommand{\braket}[1]{\langle#1\rangle}
\newcommand{\D}{\mathcal{D}}
\usepackage{cancel}
\usepackage{bm}
\usepackage{fancyhdr}
\usepackage[utf8x]{inputenc}
\usepackage[T1]{fontenc}
\usepackage[margin=0.8in,top=1in,bottom=1in]{geometry}
%%%%%
\begin{filecontents*}{refs.bib}
@article{cohen1986quantum,
  title={Quantum mechanics, volume 1},
  author={Cohen-Tannoudji, Claude and Diu, Bernard and Laloe, Frank},
  journal={Quantum Mechanics},
  volume={1},
  pages={898},
  year={1986}
}
\end{filecontents*}
%
\newcommand{\institution}{University of Arizona}
\newcommand{\autor}{Nicolás Hernández Alegría}
\newcommand{\course}{OPTI 544 Quantum Optics}
\newcommand{\assignment}{Assignment 4}
%
\title{\textbf{\assignment}\\\course\\{\Large\institution}}
\author{\autor}
%\date{\today\\Total time: 12 hours}
%
\renewcommand{\sectionmark}[1]{\markright{#1}}
\fancypagestyle{mainstyle}{
    \fancyhf{} % Clear all header and footer fields
    \fancyfoot[C]{\thepage}
    \fancyhead[LE,RO]{\course} % Section name on odd pages
    \fancyhead[LO,RE]{\assignment}
    % Optional: Thin rules
    \renewcommand{\headrulewidth}{0pt} % Header rule
    \renewcommand{\footrulewidth}{0pt} % No footer rule
}
%
\begin{document}

\pagestyle{mainstyle}
\maketitle
%%
\section*{Exercise 1}
\begin{enumerate}[itemsep=0pt,topsep=0pt,parsep=0pt,label=\alph*)]
  \item State just after the slits, and just after the linear polarizers are:
  \begin{align*}
    \ket{\Psi^-}&=\frac{1}{\sqrt{2}}(\ket{A}+\ket{B})\otimes\ket{Y},\\
    \ket{\Psi^+}&=\frac{1}{\sqrt{2}}(\ket{A}\hP_A+\ket{B}\hP_B)\otimes\ket{Y}=\frac{\sin\theta}{\sqrt{2}}(\ket{A}
    \ket{\theta_A}-\ket{B}\ket{\theta_B})\longrightarrow\ket{\Psi^{+'}}=\frac{\ket{\Psi^+}}{\sin\theta}.
  \end{align*}
  \item At the screen, $\ket{A}$ and $\ket{B}$ should evolve as we seen in the lectures. However, I will leave them as they are and make the projection.
  \begin{align*}
    \braket{x_k|\Psi^{+'}}=\frac{1}{\sqrt{2}}(\braket{x_k|A}\ket{\theta_A}-\braket{x_k|B}\ket{\theta_B})=\frac{1}{\sqrt{2}}[\psi_A(x_k)\ket{\theta_A}+
    \psi_B(x_k)\ket{\theta_B}].
  \end{align*}
  Then, the magnitude is:
  \begin{align*}
    |\braket{x_k|\Psi^{+'}}|^2&=\frac{1}{2}[\psi_A(x_k)\ket{\theta_A}+
    \psi_B(x_k)\ket{\theta_B}][\Psi^*_A(x_k)\bra{\theta_A}+
    \Psi^*_B(x_k)\bra{\theta_B}]\\
    &=\frac{1}{2}\left[|\psi_A(x_k)|^2+|\psi_B(x_k)|^2-\psi_A(x_k)\psi_B\braket{\theta_B|\theta_A}-\psi_B(x_k)\psi_A^*\braket{\theta_A|\theta_B}\right].
  \end{align*}
  We need to find the projection of the polarizers:
  \begin{align*}
    \braket{\theta_A|\theta_B}&=[\cos\theta\bra{X}+\sin\theta\bra{Y}][\cos\theta\ket{X}-\sin\theta\ket{Y}]=\cos^2\theta-\sin^2\theta=\cos2\theta,\\
    \braket{\theta_B|\theta_A}&=[\cos\theta\bra{X}-\sin\theta\bra{Y}][\cos\theta\ket{X}+\sin\theta\ket{Y}]=\cos^2\theta-\sin^2\theta=\cos2\theta.
  \end{align*}
  If we assume that 
  \begin{align*}
    \psi_A(x_k)=|\psi_A|\angle\theta_A,\quad\text{and}\quad\psi_B(x_k)=|\psi_B|\angle\theta_B,
  \end{align*}
  then usinging it with the projection computed we finally have:
  \begin{align*}
    P_k=|\braket{x_k|\Psi^{+'}}|^2=\frac{1}{2}\left[|\psi_A|^2+|\psi_B|^2-2\cos2\theta|\psi_A||\psi_B|\cos(\Delta\phi)\right],\quad\Delta\phi=\theta_B-\theta_A.
  \end{align*}
  \item For each angle, we have:
  \begin{align*}
    \theta=0:&\qquad P_k=\frac{1}{2}\left[|\psi_A|^2+|\psi_B|^2-2|\psi_A||\psi_B|\cos(\Delta\phi)\right]\Longrightarrow\text{Interference term maximized},\\
    \theta=\pi/4:&\qquad P_k=\frac{1}{2}\left[|\psi_A|^2+|\psi_B|^2\right]\Longrightarrow\text{No interference term}.
  \end{align*}
  As we can see, when $\theta=\pi/4$ the interference term vanished and therefore there is no interference.
  \item The eraser projects onto its basis:
  \begin{align*}
    \ket{\Psi^{++}}&=\ket{\theta_E}\bra{\theta_E}\ket{\Psi^{+'}}\\
    &=\frac{1}{\sqrt{2}}\ket{\theta_E}\bra{\theta_E}(\ket{A}\ket{\theta_A}-\ket{B}\ket{\theta_B})\\
    &=\frac{1}{\sqrt{2}}\left[(\cos\theta\cos\theta_E+\sin\theta\sin\theta_E)\ket{A}-(\cos\theta\cos\theta_E-\sin\theta\sin\theta_E)\ket{B}\right]\otimes\ket{\theta_E}\\
    \ket{\Psi^{++}}&=\frac{1}{\sqrt{2}}\left[\cos(\theta-\theta_E)\ket{A}-\cos(\theta+\theta_E)\ket{B}\right]\otimes\ket{\theta_E}.
  \end{align*}
  \item For $\theta=\pi/4$ and $\theta_E=\pi/4$ the state vector is:
  \begin{align*}
    \ket{\Psi^{++}}=\frac{1}{\sqrt{2}}\ket{A}\otimes\ket{\theta_E}\longrightarrow\ket{\Psi^{++'}}=\ket{A}\otimes\ket{\theta_E}.
  \end{align*}
  Because there is not interaction with the slit B, there is no interference pattern detected on the screen.
  \item For $\theta=\pi/4$ and $\theta_E=0$ we have:
  \begin{align*}
    \ket{\Psi^{++}}=\frac{\cos(\pi/4)}{\sqrt{2}}[\ket{A}-\ket{B}]\otimes\ket{X}\longrightarrow\ket{\Psi^{++'}}=\frac{1}{\sqrt{2}}[\ket{A}-\ket{B}]\otimes\ket{X}.
  \end{align*} 
  As both slits are present in the state vector, when computed the probability $|\braket{x_j|\Psi^{++'}}$ there will be a cross-term that will implies interfence.
\end{enumerate}


\section*{Exercise 2}
The intensity operator for each detector and the difference is:
\begin{align*}
  \hI_1&=\hb_1^\dagger\hb_1=\frac{1}{2}(\ha_1^\dagger+\ha^\dagger_2)(\ha_1+\ha_2)=\frac{1}{2}(\ha_1^\dagger\ha_1+\ha_1^\dagger\ha_2+\ha_2^\dagger\ha_1+\ha_2^\dagger\ha_2),\\
  \hI_2&=\hb_2^\dagger\hb_2=\frac{1}{2}(\ha_1^\dagger-\ha^\dagger_2)(\ha_1-\ha_2)=\frac{1}{2}(\ha_1^\dagger\ha_1-\ha_1^\dagger\ha_2-\ha_2^\dagger\ha_1+\ha_2^\dagger\ha_2),\\
  \hI_1-\hI_2&=\ha_1^\dagger\ha_2+\ha_2^\dagger\ha_1.
\end{align*}
We use these to compute the expected number of photons $\braket{\hI_k}$ at individual detectors and the difference $\braket{\hI_1-\hI_2}$.
\begin{enumerate}[itemsep=0pt,topsep=0pt,parsep=0pt,label=\roman*)]
  \item 
  \begin{align*}
    \braket{\hI_1}&=\frac{1}{2}\braket{1,1|(\ha_1^\dagger\ha_1+\ha_1^\dagger\ha_2+\ha_2^\dagger\ha_1+\ha_2^\dagger\ha_2)|1,1}=1,\\
    \braket{\hI_2}&=\frac{1}{2}\braket{1,1|(\ha_1^\dagger\ha_1-\ha_1^\dagger\ha_2-\ha_2^\dagger\ha_1+\ha_2^\dagger\ha_2)|1,1}=1,\\
    \braket{\hI_1-\hI_2}&=0.
  \end{align*}
  \item In this case of entangled photons, we follow the definition of the expected value and we use distribution property:
  {\small
  \begin{align*}
    \braket{\hI_1}=\left[\frac{1}{\sqrt{2}}(\bra{0,1}+\bra{1,0})\right]|\hb_1^\dagger\hb_1|\left[\frac{1}{\sqrt{2}}(\ket{1,0}+\ket{0,1})\right]=\frac{1}{4}(\bra{0,1}+\bra{1,0})|(\ha_1^\dagger\ha_1+\ha_1^\dagger\ha_2+\ha_2^\dagger\ha_1+\ha_2^\dagger\ha_2)|(\ket{1,0}+\ket{0,1}).
  \end{align*}}
  I should distribute and solve the 16 terms. However, I can express the input in the following way:
  \begin{align*}
    \ket{\psi}=\frac{1}{\sqrt{2}}(\ha_1^\dagger+\ha_2^\dagger)\ket{0,0}=\hb_1^\dagger\ket{0,0}=\ket{1',0'}.
  \end{align*}
  which is simpler to use, because:
  \begin{align*}
    \braket{\hI_1}=\braket{0',1'|\hb_1^\dagger\hb_1|1',0'}=\braket{0',1'|1',0'}=1.
  \end{align*}
  For the detector 2, the expected number of photon is:
  \begin{align*}
    \braket{\hI_2}=\braket{0',1'|\hb_2^\dagger\hb_2|1',0'}=0.
  \end{align*}
  Therefore,
  \begin{align*}
    \braket{\hI_1-\hI_2}=\braket{\hI_1}-\braket{\hI_2}=1.
  \end{align*}
  \item Using the form from previous part, the entangled state is:
  \begin{align*}
    \ket{\psi}=\frac{1}{\sqrt{2}}(\ket{1,0}-\ket{0,1})=\frac{1}{\sqrt{2}}(\ha_1^\dagger-\ha_2^\dagger)\ket{0,0}=\hb_2^\dagger\ket{0,0}=\ket{0',1'}.
  \end{align*}
  Now, 
  \begin{align*}
    \braket{\hI_1}&=\braket{0',1'|\hb_1^\dagger\hb_2'|1',0'}=0,\\
    \braket{\hI_2}&=\braket{0',1'|\hb_2^\dagger\hb_2|0',1'}=1,\\
    \braket{\hI_1-\hI_2}&=-1.
  \end{align*}
\end{enumerate}

\section*{Exercise 3}
\begin{enumerate}[itemsep=0pt,topsep=0pt,parsep=0pt,label=\alph*)]
  \item We know that the fine-structure constant is:
  \begin{align*}
    \alpha=\frac{e^2}{4\pi\varepsilon_0\hbar c}\approx\frac{1}{137}.
  \end{align*}
  So we can use it to express the ratio $r_0/\lambda_0$ in terms of $\alpha$.
  First, centrifugal force of the electron must be equal to the electrostatic force:
  \begin{align*}
    \frac{m_ev_n^2}{r_n}=\frac{1}{4\pi\varepsilon_0}\frac{e^2}{r_n^2}\longrightarrow m_ev_n^2r_n=\frac{e^2}{4\pi\varepsilon_0}.
  \end{align*}
  For $n=1$, the angular momentum $m_ev_1r_1$ must be quantized by discrete values of $\hbar$, that is,
  \begin{align*}
    m_ev_1r_1=1\hbar.
  \end{align*}
  Division of the force equation and the angular momentum yields:
  \begin{align*}
    v_1=\frac{e^2}{4\pi\varepsilon_0\hbar}
  \end{align*}
  If we substitute $v_1$ in the angular momentun equation and solve for $r_1$ we obtain:
  \begin{align*}
    r_1=\frac{\hbar}{m_ev_1}=\frac{4\pi\varepsilon_0\hbar^2}{m_ee^2}.
  \end{align*}
  We are given the resonance frequency 
  \begin{align*}
    \lambda_0=\frac{2\pi\hbar c}{E_0}=\frac{4\pi\hbar c}{m_ev_1^2},\quad E_0=\frac{1}{2}m_ev_1^2.
  \end{align*}
  Making the division $r_0/\lambda_0$, with $r_0=r_1$ in this context, and using $v_1$ yields
  \begin{align*}
    \frac{r_0}{\lambda_0}=\frac{\hbar}{m_ev_1}\frac{m_ev_1^2}{4\pi\hbar c}=\frac{v_1}{4\pi c}=\frac{1}{4\pi}\frac{e^2}{4\pi\varepsilon_0\hbar c}=\frac{\alpha}{4\pi}.
  \end{align*}
  \item The ratio $v_0/c$, with $v_0=v_1$, is exactly $\alpha$ from the velocity of last part:
  \begin{align*}
    \frac{v_0}{c}=\frac{1}{c}\frac{e^2}{4\pi\varepsilon_0\hbar}=\alpha.
  \end{align*}
  \item We will use the proportional relations:
  \begin{align*}
    \frac{r_0}{\lambda_0}\sim\alpha,\quad\omega_0\sim\frac{c}{\lambda_0}.
  \end{align*}
  {\small
  \begin{align*}
    \frac{g}{\omega_0}\sim&\frac{er_0}{\hbar\omega_0}\frac{\sqrt{\hbar\omega_0}}{\sqrt{2\varepsilon_0V}}\sim\frac{er_0V^{-1/2}}{\sqrt{\hbar\omega_0\varepsilon_0}}\sim\frac{e}{\sqrt{\varepsilon_0\hbar c}}r_0V^{-1/2}\lambda_0^{1/2}\sim\alpha^{1/2}(\alpha\lambda_0)V^{-1/2}\lambda_0^{1/2}\sim\alpha^{3/2}\lambda_0^{3/2}V^{-1/2}=\frac{\alpha^{3/2}}{(V/\lambda_0^3)^{1/2}}.
  \end{align*}}
\end{enumerate}

%\clearpage
%\nocite{*}
%\bibliographystyle{plain}   % or unsrt, alpha, apalike, etc.
%\bibliography{refs}

\end{document}
